\section{Turn complete pictures into a complete clip}
A complete picture in memory is not yet a complete output file. The clip driver
owns the input descriptor, the encoded-stream descriptor and a reconstruction
trace descriptor. It reads one complete raw picture, invokes the picture entry,
writes both results and only then advances its accepted-picture history.
The reconstruction file is an inspection aid for this assignment; ordinary H.263
players consume the encoded stream alone.

Use two reconstruction buffers. Before the first picture there is no accepted
reference, so the picture request uses reference pointer zero and INTRA mode.
The first completed reconstruction may be placed in either buffer. After its
output writes succeed, that buffer becomes the accepted reference and the other
becomes the next private destination. Continue exchanging their roles after each
accepted picture. A failed second output must not promote the staged destination.

Do not read the complete clip into a growing array. Reuse one 38\,016-byte raw
input buffer, the two reconstructed pictures, the existing picture workspace
and one encoded-picture buffer. This bounds working memory independently of
whether the input contains two or two hundred pictures. The encoded file can
grow across pictures without changing the one-picture staging contract.

\section{Repeat short operations until a complete transfer is accepted}
A file read requests a maximum number of bytes; a positive result can be smaller.
For an exact picture load, keep a destination cursor and remaining count. Request
at most 4096 bytes, advance by the \emph{returned} positive count, subtract the
same count and repeat until the remaining count reaches zero. Zero before that
point means premature end of file; a negative result means an input error.
Neither permits encoding the partially filled raw buffer.

Use the same returned-count rule for writes. A positive short write accepts
only that prefix. Retain the unaccepted suffix and retry it. A nonpositive result
ends the operation with failure rather than spinning or claiming that all requested
bytes were written. Successful completion of the encoded write still precedes
the reconstruction write and reference promotion.

The current MARS educational write syscall reads the requested range and one
additional guard byte before calling the output service. Keep that read inside
mapped private storage. The worked helper stages at most 4096 output bytes in a
4097-byte scratch window, writes a zero guard after the staged range and requests
only the actual output count. That guard is not file content. A positive short
write advances the original source by the accepted count; the next attempt stages
the still-unaccepted suffix. This preserves bounded accesses without imposing an
extra readable byte on the caller's picture buffers.

Keep nested-call state in saved registers or stack records. A descriptor is an
integer returned by open, not a fixed constant such as three. Close every descriptor
successfully opened by this call on every return path. An output-open failure
can occur after the input has opened; cleanup must still close the input. Do not
close an unowned negative descriptor or the process's standard streams.

\section{Distinguish accepted progress from whole-clip success}
A useful failure result reports the last accepted picture, while still marking
the entire clip as failed. The progress record has three native-endian doublewords:
\begin{center}
\tagpdfsetup{table/header-rows={1}}
\begin{tabular}{rll}
\textbf{Offset} & \textbf{Field} & \textbf{Meaning}\\
0 & Accepted pictures & Both required writes completed for each\\
8 & Accepted encoded bytes & Sum for those pictures only\\
16 & Accepted reference & Pointer to the last such reconstruction; zero initially
\end{tabular}
\end{center}
Initialize this record only after complete metadata and path admission. Update
all three fields after each accepted picture. A later input error, output error
or picture-child failure retains that accepted prefix and its reference samples.
The destination buffer for the rejected picture may already contain private work.

Consider a first encoded picture of 696 bytes followed by another picture. If
the second stream write accepts ten bytes and then fails, the stream contains
706 bytes, but accepted encoded bytes remain 696. If the second stream write
finishes and its reconstruction write then fails, the file can contain a complete
second encoded picture that is nevertheless unaccepted by this driver. Its
progress must still name the first reconstruction.

After all requested pictures, read one extra input byte to require exact EOF.
A positive result reveals trailing input; a negative result is still an input
error. These cases can occur after every picture output has been accepted.
The progress count then equals the requested count, but overall status is failure.
A nonzero status means the run is not a finished clip, regardless of how much of
its output prefix can be decoded. Files already written are not rolled back.

\section{Pass an explicit clip request}
Implement \texttt{h263\_encode\_clip} with an eight-byte-aligned 104-byte request
in \texttt{\$a0} and an eight-byte-aligned 24-byte progress result in
\texttt{\$a1}. Fields are unsigned native-endian doublewords:
\begin{center}
\tagpdfsetup{table/header-rows={1}}
\begin{tabular}{rll}
\textbf{Offset} & \textbf{Field} & \textbf{Meaning}\\
0, 8, 16 & Input, stream, reconstruction paths & Distinct nonempty NUL-terminated names\\
24 & Pictures & One through 299\\
32 & Clock step & One through 255 baseline ticks\\
40 & Quantizer & One through 31\\
48 & Arithmetic mode & Zero scalar; one MSA\\
56 & Raw input buffer & 38\,016 active bytes\\
64, 72 & Reconstruction A, B & 38\,016 active bytes each\\
80 & Picture workspace & Sixteen-byte-aligned 169\,088-byte arena\\
88 & Encoded-picture buffer & 131\,072 active bytes\\
96 & Reserved & Zero
\end{tabular}
\end{center}
Require $N\times T\times1001\leq300\,000$, where $N$ is picture count and $T$
is clock step. Validate each bounded factor before multiplying. This is the
same ten-second limit used by the generator, including the final picture hold.
The baseline clock never exceeds thirty pictures per second.

Each path supplies a mapped 4096-byte window with its first NUL no later than
byte 4095. Validate all record/data extents and controls before scanning paths;
then reject empty, unterminated or equal path names before opening anything.
Compare path contents, not just their pointers. Do not truncate an input by opening
the same name for output through a different string buffer. The preparation helper
uses printable ASCII names for the current byte-oriented syscall string ABI.

All required extents are mapped and non-wrapping. Request, progress, paths,
buffers, workspace and private stack are disjoint. Preserve request/path bytes,
saved GPRs and the incoming sixteen-byte-aligned stack pointer. Supply 65\,536
private stack bytes. Raw, reconstruction, workspace and encoded buffers are
working storage and may change after valid admission, including on failure.
The accepted reference named by progress must still contain its last accepted
picture after a later failure.

Return status zero and accepted-picture count on success. Invalid arguments
return one/zero without changing progress, opening files or touching data buffers.
For admitted calls, the following statuses return the actual accepted count and
retain the corresponding progress record:
\begin{center}
\tagpdfsetup{table/header-rows={1}}
\begin{tabular}{rl}
\textbf{Status} & \textbf{Meaning}\\
2 & A required file could not be opened\\
3 & Read error, incomplete picture or trailing input\\
4 & Picture child failed or returned an invalid byte count\\
5 & A required output write failed
\end{tabular}
\end{center}
The supplied unfinished clip entry uses development status seven. The nested
picture entry has its own status meanings; the driver maps any failed child to
its clip-level status four. All opened descriptors are closed before returning.
